\ifdefined\PortfolioMain
\else
\documentclass[11pt,a4paper]{article}

% Eigenständiges individuelles Logbuch für das Santander-Projekt.
\usepackage{fontspec}
\setmainfont{Latin Modern Roman}
\usepackage[a4paper,top=2.5cm,bottom=2.5cm,left=2.7cm,right=2.7cm]{geometry}
\usepackage{xcolor}
\definecolor{ADAblue}{HTML}{173F5F}
\definecolor{ADAgold}{HTML}{D99A2B}
\usepackage{booktabs}
\usepackage{longtable}
\usepackage{array}
\usepackage{tabularx}
\usepackage{hyperref}
\hypersetup{
  colorlinks=true,
  linkcolor=ADAblue,
  urlcolor=ADAblue,
  pdftitle={Logbuch Santander Customer Transaction Prediction},
  pdfauthor={Cyprian Sawarzynski}
}

\newcommand{\placeholder}[1]{\textcolor{ADAgold}{\textbf{[PLACEHOLDER: #1]}}}
\newcommand{\logsection}[2]{\subsection*{#1 -- #2}}
\newcommand{\field}[2]{\noindent\textbf{#1}\\#2\par\medskip}

\title{\color{ADAblue}\textbf{Individuelles Logbuch}\\[0.4em]
       \large Santander Customer Transaction Prediction}
\author{Cyprian Sawarzynski\\Gruppe 138}
\date{Sommersemester 2026}
\fi

\providecommand{\placeholder}[1]{\textcolor{ADAgold}{\textbf{[PLACEHOLDER: #1]}}}
\providecommand{\logsection}[2]{\subsection*{#1 -- #2}}
\providecommand{\field}[2]{\noindent\textbf{#1}\\#2\par\medskip}

\ifdefined\PortfolioMain
\else
\begin{document}
\maketitle
\tableofcontents
\clearpage
\fi

\section{Zweck und Führungsregeln}

Dieses Logbuch dokumentiert meinen individuellen Arbeitsprozess während des gesamten
Bearbeitungszeitraums. Es ist kein nachträglich geglätteter Ergebnisbericht, sondern hält
den tatsächlichen Weg zur Analyse fest. Dadurch sollen Arbeitsaufwand, eigene Entscheidungen,
Anpassungen, Schwierigkeiten und verworfene Ansätze nachvollziehbar bleiben.

Die Einträge werden chronologisch und möglichst unmittelbar nach dem jeweiligen Arbeitsschritt
ergänzt. Jeder Eintrag enthält mindestens:

\begin{itemize}
  \item Datum des Eintrags,
  \item einen kurzen und prägnanten Titel,
  \item eine verständliche Zusammenfassung des Arbeitsschritts.
\end{itemize}

Zusätzlich dokumentiere ich systematisch:

\begin{itemize}
  \item konkrete Arbeitsschritte und Ergebnisse,
  \item methodische Entscheidungen und deren Begründung,
  \item verworfene Ansätze und Iterationen,
  \item Schwierigkeiten, Unsicherheiten und offene Fragen,
  \item Zeitaufwand und gegebenenfalls Bezug zu Gruppentreffen,
  \item nächste Schritte und Abweichungen vom vereinbarten Vorgehen.
\end{itemize}

Das Logbuch bildet die Grundlage für die spätere Zusammenfassung im E-Portfolio. Abbildungen
oder Tabellen dürfen hier auch vorläufig bzw. als Arbeitsdokumentation eingefügt werden;
für die finale Zusammenfassung werden die relevanten Darstellungen später formal ausgearbeitet.

\section{Projekt- und Gruppenrahmen}

\begin{tabularx}{\textwidth}{>{\bfseries}p{0.32\textwidth}X}
  Thema & Santander Customer Transaction Prediction \\
  Datensatz & OpenML ID 45566 \\
  Name & Cyprian Sawarzynski \\
  Gruppennummer & 138 \\
  Gruppenmitglieder & Tim Hebestreit (Gruppenleiter), Angeline Usanayo, Trong Cao \\
  Sprecher/Gruppenleitung & Tim Hebestreit \\
  Projektdeadline & 01.09.2026 \\
  Individueller Workload-Richtwert & ca. 100 Stunden \\
\end{tabularx}

\section{Chronologische Einträge}

\logsection{06.08.2026}{Projektunterlagen und Projektauswahl}
\field{Zusammenfassung}{Die Bewertungs- und Projektunterlagen wurden gelesen. Als Projekt aus
den vorhandenen Aufgaben wurde Santander Customer Transaction Prediction ausgewählt.}
\field{Arbeitsschritte}{
\begin{itemize}
  \item Bewertungshinweise für das E-Portfolio gelesen.
  \item Anforderungen an EDA und Machine Learning herausgearbeitet.
  \item Santander-Aufgabenblatt gelesen und OpenML-Datensatz ID 45566 notiert.
  \item Kaggle- und OpenML-Links als spätere Quellen vorgemerkt.
\end{itemize}}
\field{Entscheidungen und Begründung}{Die Analyse wird auf Santander fokussiert, weil die Aufgabe eine
binäre Transaktionsvorhersage mit klarer OpenML-Referenz vorgibt. Die weitere Bearbeitung soll
EDA-Erkenntnisse ausdrücklich mit den späteren Modellentscheidungen verbinden.}
\field{Schwierigkeiten und offene Fragen}{Die genaue Aufgabenverteilung innerhalb der Gruppe,
die persönliche Forschungsfrage, die Modellstrategie und der verfügbare Rechenrahmen sind noch
zu klären.}
\field{Abweichung vom Vorgehen}{Noch keine Abweichung vom Gruppenplan dokumentiert.}
\field{Gruppentreffen}{\placeholder{kein Treffen / Datum und Kurzverweis ergänzen}}
\field{Zeitaufwand}{\placeholder{Minuten/Stunden ergänzen}}
\field{Nächster Schritt}{Logbuch und E-Portfolio als belastbare Arbeitsgrundlage anlegen und anschließend
den Datensatz laden sowie technisch prüfen.}

\logsection{06.08.2026}{E-Portfolio auf LaTeX umgestellt}
\field{Zusammenfassung}{Das zunächst als Notebook angelegte E-Portfolio wurde auf ein eigenständiges
LaTeX-Dokument umgestellt. Ziel war ein besser kontrollierbares Layout für die spätere PDF-Abgabe.}
\field{Arbeitsschritte}{
\begin{itemize}
  \item Ordner \texttt{E-Protfolio} im eigenen Projektordner angelegt.
  \item Hauptdatei \texttt{ADA\_138\_Sawarzynski\_Portfolio.tex} erstellt.
  \item Titel, Projektüberblick, Logbuch-Platzhalter, Zusammenfassung, Gruppenreflexion,
        Quellen und Abschluss-Checkliste strukturiert.
  \item Ordner \texttt{assets} für spätere Abbildungen vorbereitet.
  \item LaTeX-Datei mit XeLaTeX kompiliert und PDF-Ausgabe geprüft.
\end{itemize}}
\field{Entscheidungen und Begründung}{Das Logbuch wird als separate Datei geführt, damit der
Arbeitsprozess unabhängig vom später verdichteten Portfolio fortlaufend ergänzt werden kann.}
\field{Schwierigkeiten und Lösungen}{Die lokale TinyTeX-Installation enthielt mehrere optionale
Pakete nicht. Die Präambel wurde deshalb auf vorhandene Standardpakete reduziert; die Ausgabe
kompiliert nun erfolgreich.}
\field{Abweichung vom Vorgehen}{Das ursprünglich geplante Notebook-Format wurde zugunsten von
LaTeX verworfen, nachdem die PDF-Gestaltung von nbconvert nicht zufriedenstellend war.}
\field{Gruppentreffen}{Kein inhaltliches Gruppentreffen dokumentiert.}
\field{Zeitaufwand}{\placeholder{Minuten/Stunden ergänzen}}
\field{Nächster Schritt}{Die EDA des Santander-Datensatzes beginnen und dabei Datenform, Datentypen,
Zielvariable, fehlende Werte, Klassenverteilung und erste Verteilungen systematisch untersuchen.}

\logsection{\placeholder{Datum}}{\placeholder{prägnanter Arbeitstitel}}
\field{Zusammenfassung}{\placeholder{Kurze Zusammenfassung des konkreten Arbeitsschritts.}}
\field{Arbeitsschritte}{\placeholder{Welche Schritte wurden tatsächlich durchgeführt?}}
\field{Ergebnis/Erkenntnis}{\placeholder{Was kam dabei heraus? Zahlen, Tabellen oder Abbildungen verlinken.}}
\field{Entscheidung und Begründung}{\placeholder{Welche Entscheidung wurde getroffen und warum?}}
\field{Verworfene Ansätze/Schwierigkeiten}{\placeholder{Was hat nicht funktioniert oder wurde verworfen?}}
\field{Abweichung vom vereinbarten Vorgehen}{\placeholder{Keine / konkrete Abweichung mit Begründung.}}
\field{Gruppentreffen}{\placeholder{Datum, Teilnehmende und Bezug zum Treffen.}}
\field{Zeitaufwand}{\placeholder{Minuten/Stunden.}}
\field{Nächster Schritt}{\placeholder{Konkreter nächster Arbeitsschritt.}}

\section{Laufende Methodik- und Entscheidungsübersicht}

Diese Übersicht wird parallel zum Logbuch aktualisiert. Sie verhindert, dass wichtige
Entscheidungen nur im Fließtext verborgen bleiben.

\begin{longtable}{p{0.19\textwidth}p{0.29\textwidth}p{0.42\textwidth}}
\toprule
\textbf{Bereich} & \textbf{Aktueller Stand} & \textbf{Begründung/Beleg} \\
\midrule
Datenzugang & \placeholder{offen} & \placeholder{Quelle, Version, Zugriffsdatum} \\
Datentypen & \placeholder{offen} & \placeholder{Speicherbedarf und geeignete Wahl} \\
Fehlende Werte & \placeholder{offen} & \placeholder{EDA-Befund und Umgang damit} \\
Zielvariable & \placeholder{offen} & \placeholder{Kodierung und Klassenverteilung} \\
Train-Test-Aufteilung & \placeholder{offen} & \placeholder{Zeitpunkt und Begründung} \\
Metrik(en) & \placeholder{offen} & \placeholder{Warum passend für die Aufgabe?} \\
Modelle & \placeholder{offen} & \placeholder{Vergleich und Auswahl} \\
Feature Engineering & \placeholder{offen} & \placeholder{getestete Varianten} \\
Hyperparameter & \placeholder{offen} & \placeholder{Suchraum und Evaluation} \\
\bottomrule
\end{longtable}

\section{Offene Fragen und Risiken}

\begin{itemize}
  \item \placeholder{Welche genaue persönliche Forschungsfrage wird verfolgt?}
  \item \placeholder{Wie wird die Gruppenarbeit aufgeteilt?}
  \item \placeholder{Welche Metrik ist für die unausgeglichene Zielvariable sachgerecht?}
  \item \placeholder{Wie groß ist der Datensatz und welche Speicherstrategie ist erforderlich?}
  \item \placeholder{Welche Grenzen der Aussagekraft müssen später diskutiert werden?}
\end{itemize}

\section{Hinweise für die spätere Zusammenfassung}

Bei der späteren Verdichtung des Logbuchs müssen Fragestellung, Methodenwahl, Ergebnisse,
Schwierigkeiten/Grenzen und Ausblick klar erkennbar sein. Für Abbildungen und Tabellen werden
Nummer, kurze Caption, Textreferenz und ein prägnanter Take-away ergänzt. Die ausführliche
Nachvollziehbarkeit bleibt in diesem Logbuch erhalten.

\ifdefined\PortfolioMain
  \def\LogbuchEnd{}
\else
  \def\LogbuchEnd{\end{document}}
\fi
\LogbuchEnd
